\def\stat{fedoseev}





\def\tit{ЭЛЕКТРОННЫЕ ОБРАЗОВАТЕЛЬНЫЕ РЕСУРСЫ: 


ЭФФЕКТИВНОСТЬ ПРИМЕНЕНИЯ}





\def\titkol{Электронные образовательные ресурсы: 


эффективность применения} 





\def\autkol{А.\,А.~Федосеев}





\def\aut{А.\,А.~Федосеев$^1$}





\titel{\tit}{\aut}{\autkol}{\titkol}





%{\renewcommand{\thefootnote}{\fnsymbol{footnote}}\footnotetext[1]


%{Работа


%выполнена при поддержке Российского фонда фундаментальных


%исследований, проекты 11-01-12026-офи-м, 10-07-00021.}}





\renewcommand{\thefootnote}{\arabic{footnote}}


\footnotetext[1]{Институт проблем информатики Российской академии наук, A.Fedoseev@ipiran.ru}








\Abst{На основе анализа процессов научения, учения и обучения выделено 


информационное содержание  учебного процесса и построена его информационная 


модель. Определено место в учебном процессе желательного приложения усилий по 


созданию электронных образовательных ресурсов (ЭОР). Определены желательные свойства 


ресурсов для обеспечения их положительного влияния на эффективность учебного 


процесса.}





\KW{учение; обучение; мотивация; восприятие; закрепление; учебный процесс; 


информационная модель; электронный образовательный ресурс}





 \vskip 14pt plus 9pt minus 6pt








      \thispagestyle{headings}





      


            \label{st\stat}


            


            \vspace*{-12pt}





\section{Введение}





В области применения информационно-коммуникационных технологий (ИКТ) в 


образовании сложилась следующая картина: на протяжении десятилетий в эту область 


вкладываются серьезные средства, школы оснащены современной вычислительной 


техникой, подключены к Интернету. Существуют и доступны для применения в школах 


или в индивидуальном порядке десятки и сотни тысяч ЭОР. 


Каждый год к ним прибавляются тысячи новых. Опережающими темпами 


создаются методические работы по их применению. Учителя активно обмениваются 


опытом по использованию ЭОР. Параллельно развиваются и совершенствуются 


комплексы требований, которым должны удовлетворять создаваемые ЭОР, в том числе по 


критериям актуальности содержания, дидактической проработанности, научности, 


принимаются стандарты, разрабатываются механизмы обеспечения качества ЭОР. Тем не 


менее всего этого оказывается недостаточно: существенного положительного эффекта от 


использования средств ИКТ не наблюдается даже в простейшем случае~--- при 


использовании критерия лучшего освоения материала учащимися за то же время или 


такого же уровня освоения за меньшее время. Надо полагать, что при обнаружении 


заметного влияния применения ЭОР на эффективность учебного процесса в школах, об 


этом немедленно стало бы известно. Но таких публикаций пока не появилось. Напротив, 


отсутствие эффекта отмечается как у нас в стране, так и за рубе-\linebreak жом~[1, 2].





В статье сделана попытка определить, следует ли вообще ожидать положительного 


эффекта от существующих и продолжающих множиться в огромных количествах ЭОР. 


Также показано, какими свойствами должны обладать ЭОР, чтобы можно было ожидать 


эффекта от их применения.








\section{Теоретические основы}





Педагогическая психология~[3] рассматривает приобретаемые в результате учебного 


процесса знания, умения и навыки учащихся как их индивидуальный опыт. Способность к 


приобретению индивидуального опыта является общим свойством биологических систем. 


Процесс приобретения опыта называется научением. Человеку, в отличие от других видов 


биологических систем, свойственно целенаправленно приобретать свой индивидуальный 


опыт. Этот процесс, в отличие от бессознательного научения, носит название учения. На 


рис.~1~[4] перечислены виды деятельности, в результате которых приобретается опыт. 


Учение охватывает те <<деятельности>>, которые находятся в столбце <<преднамеренное 


познание>>. Все остальные ветви схемы отображают процесс научения.





Роль учителя в процессах учения и научения (для учащихся младшего возраста) 


заключается в организации передачи социокультурного опыта учащимся и тех видов 


деятельности, которые позволят учащимся приобрести этот опыт уже в виде своего 


индивидуального опыта. В~этих процессах учитель занимается обучением. Таким 


образом, учитель организует передачу опыта учащимся и в зависимости от их возрастных, 


интеллектуальных и психологических особенностей выбирает наиболее подходящие для 


этого виды деятельности. Обратим внимание, что среди видов деятельности, приводящих 


к присвоению опыта, особенно выделены такие виды, как труд, игра и стихийное 


общение, а также стихийное наблюдение, чтение, просмотр фильмов и~др. Здесь важно, 


что, во-первых, научение есть результат активности самого учащегося и, во-вторых, 


научение происходит в результате деятельности учащегося. Таким образом, какие бы 


действия ни предпринимал учитель, пока он не организовал, не запустил, не включил 


собственную активность и деятельность учащегося, никакого процесса научения 


произойти не может. И~тем более не может произойти процесса учения как осознанной 


деятельности.








Когда мы говорим, что учитель обучает, мы имеем в виду, что он передает учащимся 


учебную информацию, задает им вопросы, чтобы убедиться, что информация воспринята, 


задает задания, чтобы закрепить полученные учащимися знания. Но все это возможно 


только тогда, когда сопровождается собственной деятельностью учащихся, направленной 


на приобретение нового индивидуального опыта, соответствующего информации, 


переданной учителем. Поскольку передавать информацию, а также задавать задания 


учитель может не непосредственно, а с использованием средств ИКТ, смысл его 


деятельности перемещается на обеспечение познавательной активности учащихся. Если 


эта активность не включена, грубо говоря, учащийся витает в облаках и не собирается 


ничего познавать, то учитель не в состоянии ничему такого ученика научить. Процесс 


учения может идти только на основе собственной деятельности учащегося и никак иначе.





\begin{figure} %fig1


  \vspace*{1pt}


\begin{center}


\mbox{%


\epsfxsize=125.394mm


\epsfbox{fed-1.eps}


}


\end{center}


\vspace*{-14pt}


\Caption{Виды деятельности, в результате которых человек приобретает опыт}


\vspace*{-6pt}


\end{figure}








Согласно~[5] <<принципиальная логика современного обучения состоит в том, что она 


неизбежно начинается с предъявления учащимся разными способами (обязательно с 


помощью материальных предметов или знаков, вербально и~т.\,д.) новой информации, 


продолжается организованным репродуцированием (в разных вариантах по видам и 


сложности) информации и связанных с ней действий>>. Следовательно, если говорить о 


структуре деятельности учащегося в процессе учения (научения), то он, во-пер\-вых, 


воспринимает предъявляемое и, во-вторых, воспроизводит или копирует (репродуцирует) 


его. При этом воспроизведение является учебной деятельностью по закреплению 


предъявляемого материала: как сказано в~[6], <<закрепление и совершенствование знаний 


и привитие умений и навыков>>. Изложенное позволяет сформулировать два элемента 


структуры учебной деятельности: восприятие и закрепление.





Учитывая упомянутую выше необходимость собственной деятельности для реализации 


процесса учения, уже выделенных элементов~--- восприятия и закрепления~--- 


недостаточно. В~самом деле, далее в~[6] сказано: <<роль педагога заключается в том, 


чтобы вызвать у школьников собственную потребность овладеть новым понятием или 


законом науки, способами применения их на практике>>. Таким образом, в дополнение к 


выделенным элементам структуры учебной деятельности следует добавить особый вид 


деятельности: выработку учащимся потребности к учению. Другими словами, речь идет о 


мотивации. Можно указать работы (например,~[7]), в которых мотивация выделяется как 


самостоятельный элемент структуры учебной деятельности. Окончательно имеем: 


структура учебной деятельности учащегося представляет собой последовательность 


следующих собственных действий:


\begin{itemize}


\item мотивация к предстоящему этапу (теме) процесса учения;\\[-15pt]


\item восприятие учебного материала по теме;\\[-15pt]


\item закрепление материала по теме.


\end{itemize}





Соответствующие этим проявлениям активности учащихся действия учителя выглядят 


следующим образом:


\begin{itemize}


\item запуск мотивации к предстоящему этапу (теме) процесса учения;\\[-15pt]


\item предоставление, передача учебного материала по теме;\\[-15pt]


\item руководство процессом закрепления материала по теме.


\end{itemize}





Таким образом, роль учителя заключается в организации и обеспечении всех элементов 


процесса. Научить без соответствующей активности самого учащегося учитель ничему не 


может. Однако он может (и должен) запустить процесс мотивации учащихся, передать 


непосредственно или опосредованно новый материал и обеспечить все необходимое для 


закрепления материала.





Может показаться, что автор упустил такой важный элемент структуры учебной 


деятельности, как контроль. В~этом обвинении есть резон, поскольку, во-пер\-вых, в 


большинстве педагогических произведений необходимость контроля после закрепления 


предусмотрена и, во-вто\-рых, всем известны такие виды учебной деятельности, как 


контрольная работа, тестирование, экзамен. Тем не менее автор сознательно опускает этот 


элемент структуры по причинам, которые станут понятны позже. Здесь же будет 


уместным следующий комментарий: когда говорится об учебной деятельности, 


подразумеваются те виды деятельности, которые приводят к появлению нового 


индивидуального опыта в виде знаний, умений и навыков. Опыт появляется в результате 


восприятия и закрепления. Чтобы восприятие и закрепление состоялись, должна быть 


осуществлена мотивация собственной познавательной деятельности. Контроль не 


является видом деятельности, в результате которой появляется новый опыт. Контроль~--- 


это такое мероприятие, которое делает наличие (или отсутствие) нового опыта у учащихся 


очевидным для учителей, родителей, школы, государства. Отметим, что необходимость 


контроля возникает из ситуации, когда в учебной деятельности предполагается 


возможность (и даже неизбежность) брака, именно для обнаружения этого брака.





Рассмотрим теперь более подробно выделенные элементы структуры учебной 


деятельности.





\section{Мотивация}





У каждого человека (и учащегося в том числе) существует и действует собственная 


многофакторная система мотивации. В~зависимости от вида предпринимаемой 


деятельности преобладают различные мотивационные факторы. При необходимости эти 


факторы могут быть усилены или ослаблены специальным воздействием. Учитель имеет 


возможность хорошо узнать своих учеников, чтобы индивидуально и эффективно 


воздействовать на них. Вместе с тем существуют общие для всех закономерности. Эти 


закономерности лучше всего изучены в об\-ласти маркетинга и рекламы, где десятилетиями 


изучаются способы воздействия на человека, побуждающие его произвести определенные 


действия. Единственное отличие от коммерческой рекламы заключается в том, что 


учащийся побуждается к приобретению не товаров и услуг, а индивидуального опыта в 


виде знаний, умений и навыков. Следовательно, мотивационные материалы должны и 


могут создаваться с использованием теоретических наработок в области маркетинга и 


рекламы.





Вместе с тем педагогической наукой накоплены собственные методы мотивации 


учащихся. Одним из самых эффективных по праву считается проектный метод, 


позволяющий выстроить деятельность учащихся таким образом, чтобы достижение ими 


значимых для них целей было обусловлено приобретением определенных знаний, умений 


и навыков в соответствии с программами обучения. В~настоящее время в педагогической 


литературе существует множество подробно описанных проектов.





Как минимум учащемуся должна быть передана, а им воспринята следующая информация 


о предмете целиком или отдельной теме (группы тем) в предмете:


\begin{itemize}


\item место предмета в системе знаний и культуре или темы в предмете. Как сказал Ян 


Коменский~[8], <<при образовании юношества все нужно делать как можно более 


отчетливо, так, чтобы не только учащий, но и учащийся понимал без всякого 


затруднения, где он находится и что он делает>>;


\item значение предмета в системе знаний и культуре, в том числе значение именно для 


учащихся определенного возраста, в соответствии с их интересами, или значение темы 


или группы тем в изучаемом предмете;


\item преимущества, которыми обладает человек, овладевший предметом (темой) (и 


снова Коменский~[8]: <<Ты облегчишь ученику усвоение, если во всем, чему бы ты 


его ни учил, покажешь ему, какую это может принести повседневную пользу в 


общежитии>>).


\end{itemize}





Наличие мотивации у учащихся проявляется в возникновении вопросов, осуществлении 


попыток найти дополнительный материал или какой-либо иной активности по изучаемой 


теме (предмету). Эти действия учащихся являются сигналами учителю о возникновении 


мотивации.





\section{Восприятие}





Учащийся воспринимает ту информацию, которую непосредственно или опосредованно 


передает учитель. При этом речь может идти об осознанном восприятии, и тогда это~--- 


процесс учения, или о восприятии неосознанном, что характеризует процесс научения. 


Выбор между этим двумя процессами, а следовательно, и выбор стиля подачи и 


содержания материала осуществляет учитель, принимая во внимание возрастные, 


психологические и интеллектуальные особенности учащихся.





В тех случаях, когда требуется освоение действий: движений, поз, произнесения слов и 


выражений (в том числе на других языках), произнесения слов и выражений с упором на 


интонацию или акцент, оперирования предметами, в том числе: инструментами, 


спортивными снарядами, средствами письма и рисования, рисования линий и 


изображений, написания букв, написания слов и~т.\,п.,~--- передача информации 


заключается в демонстрации действия, в том чис\-ле, при необходимости, замедленно, в 


деталях, с обеспечением повтора необходимое чис\-ло раз. Словесные (при 


непосредственной передаче информации учителем) или текстовые комментарии являются 


в этом случае второстепенными.





В ситуациях, когда информация должна наполнить индивидуальный опыт учащихся в 


виде нового знания, в зависимости от того, кто является инициатором нового опыта~--- 


учитель или сам ученик,~--- процесс восприятия будет разным. Если информацию 


передает учитель в виде своей речи, предлагаемого текста или мультимедийного 


продукта, то учащийся воспринимает то, что передает учитель. Если же у учащегося есть 


инициативные вопросы, на которые следует найти ответы, или гипотезы, которые следует 


проверить, то никакой заранее приготовленный материал не может решить проблему. 


Учащемуся следует предоставить возможность искать ответы на свои вопросы. 


В~качестве источников знания в данном случае может выступить эрудиция учителя или 


других людей, библиотека и, разумеется, Интернет. Что касается проверки гипотез, то 


предоставление возможности для экспериментирования в виде соответствующей среды и 


времени позволяет обеспечить этот вид приобретения знаний. Далеко не всегда для 


создания среды для экспериментирования требуются серьезные затраты. Например, дети 


постоянно экспериментируют со своей речью и поведением, чтобы проверить гипотезы 


относительно допустимости произнесения некоторых слов или осуществления некоторых 


действий. Другим примером является работа с компьютером. Подавляющее большинство 


прикладных программ устроено таким образом, что позволяет экспериментировать с ними 


без фатальных повреждений с целью проверки ожидаемой реакции на не определенные 


заранее управляющие воздействия.





Необходимыми элементами этапа восприятия являются входная проверка наличия у 


учащихся знаний, умений и навыков, необходимых для восприятия нового материала, и 


обратная связь о фактическом восприятии нового материла. Очевидно, что материал не 


может быть воспринят и усвоен, если он опирается на отсутствующие у учащегося 


элементы знаний. Столь же очевидно, что выслушанные объяснения учителя не 


обязательно были услышаны, прочитанный текст мог не оставить о себе никаких 


воспоминаний, а просмотренный видеоклип не породил никаких образов и выводов в 


сознании учащегося. Только когда есть убежденность в том, что материал воспринят, его 


закрепление (следующий этап учебной деятельности) может быть эффективным.





\section{Закрепление}





Каким образом закрепляются новые знания, умения и навыки?





В случае, когда требуется закрепление действий, разумным механизмом является 


фиксация и отображение действия учащегося таким же или похожим образом, как была 


сделана демонстрация. Это позволит очевидным образом сопоставить действия учащегося 


с эталонными (демонстрируемыми) действиями. Величина отличия производимого 


действия от эталонного позволяет сделать заключение о достигнутой или нет 


правильности выполнения действия. Последующая серия повторений действия, возможно, 


в разбивку с другими действиями, с перерывами в осуществлении закрепляемого действия 


позволит убедиться в действительном его закреплении.





Новое знание, содержащееся в результатах исследований и экспериментов, закрепляется 


написанием соответствующих отчетов о самостоятельной работе, докладами и защитой в 


классе, а также, возможно, и на межшкольных мероприятиях. Сами по себе эти действия 


являются достаточными для закрепления нового знания, добытого собственными 


усилиями в ответ на сформулированные самостоятельно вопросы или созданные 


самостоятельно гипотезы.





Новое, воспринятое из материалов, переданных учителем, знание закрепляется с помощью 


упражнений и задач, опять-таки заданных учителем. Какими свойствами должен обладать 


комплект задач (упражнений), чтобы он был способен обеспечить закрепление 


воспринятого учащимся материала? Во-первых, этот комплект должен быть полным. Ни 


одна деталь переданного материала, требующая закрепления, не должна быть обойдена 


задачами и/или упражнениями. Во-вто\-рых, комплект должен содержать все типы задач, 


связанные с закрепляемым материалом и имеющие значение для дальнейшего обучения. 


В-третьих, в случаях неправильного решения задач, комплект должен содержать 


однотипные задачи с различными числовыми данными, с тем чтобы при возникновении 


необходимости повторного решения задач (например, в случаях неправильного решения) 


учащемуся не были предъявлены в точности те же задачи. У~учителя должны быть в 


распоряжении комплекты задач различного уровня сложности, чтобы снабжать учащихся 


задачами соответственно уровню их подготовки или заинтересованности. Каждый случай 


неправильного решения задачи должен сопровождаться анализом с целью выявления 


недостаточно понятых деталей материала. После прояснения материала учащемуся 


предъявляются такие же задачи, но с иными числовыми данными. Когда оказывается, что 


учащийся свободно и быстро справляется со всеми задачами комплекта, учитель может 


сделать обоснованный вывод о состоявшемся закреплении.





\section{Информационная модель учения}





В предыдущем разделе рассмотрена структура учебного процесса с точки зрения 


информатики: описаны сопровождающие учебный процесс информационные потоки. 


Оставляя в стороне вопросы, связанные с психологической сущностью процессов 


восприятия и закрепления, автору удалось сосредоточиться на информационной стороне 


дела, что позволило сформировать основную информационную модель учения в виде, 


представленном на рис.~2. При этом предполагается, что все материалы, необходимые для 


учебного процесса, у учителя имеются, что можно считать достаточно близким к 


реальности. Поэтому в предлагаемой модели учитель не занимается поиском 


соответствующего материала, а сразу передает учащимся имеющийся материал. Под 


материалом здесь имеется в виду любой носитель информации с данными. Это может 


быть речь (лекция) учителя, текст, в том числе в электронном виде, видео и прочее.





\begin{figure} %fig2


  \vspace*{1pt}


\begin{center}


\mbox{%


\epsfxsize=122.813mm


\epsfbox{fed-2.eps}


}


\end{center}


\vspace*{-12pt}


\Caption{Схема основной информационной модели учения}


\end{figure}





Процесс начинается с передачи учащимся мотивирующего материала. Восприятие этого 


материала должно ознаменоваться возникновением побуждения учащихся к овладению 


учебным материалом. Чтобы двигаться дальше, учитель должен получить обратную связь 


об уровне мотивации (на схеме~--- отчет). Фактически на уроке, видя глаза учащихся и 


ощущая их возбуждение, учитель такую обратную связь имеет, особенно когда 


мотивирующий материал заключается в выступлении самого учителя. Учитель также 


может опросить некоторых учащихся класса, чтобы более точно оценить степень 


возникшей мотивации. После чего он принимает решение о возможности двигаться 


дальше. Следует отметить, что, как правило, мотивация оценивается учителем не в виде 


специального действия, а между делом, не задерживаясь на этом аспекте учебного 


процесса специально. Возникло у учителя ощущение, что все в порядке, значит можно 


двигаться дальше~--- переходить к следующему этапу учебного процесса. Тем более речь 


не может идти об оценке уровня мотивации каждого из учащихся класса. Следовательно, 


на этом этапе учебного процесса закладывается запланированный брак: не получив 


обратной связи о достаточном уровне мотивации всех учеников, учитель переходит к 


передаче нового учебного материала. Те учащиеся, которые не обрели достаточной 


мотивации, не будут воспринимать учебный материал с должным усердием и 


результативностью. Поскольку уровень мотивации всех учащихся не оценен, а 


недостаточная мотивация не даст возможности некоторым из учащихся воспринять в 


достаточной степени учебный материал, то, следовательно, неизбежно появление 


отстающих учеников.





Как бы то ни было, учитель принимает решение передать учащимся новый учебный 


материал по теме. Здесь, как и в каждом цикле учебного процесса, требуется обратная 


информация от учащихся. Как правило, учитель получает информацию о восприятии 


учебного материала методом выборочного опроса. О~восприятии не опрошенных 


учащихся информация учителю не поступает. На этом этапе также возможен учебный 


брак: из-за появления не выявленных учителем учеников, которые не восприняли учебный 


материал в должной мере.


{\looseness=1





}





Тем не менее учитель принимает решение перейти к закреплению материала и выдает 


задания. Следует отметить, что, как правило, состав и количество заданий совершенно 


недостаточны для закрепления той порции знаний, умений и навыков, которая определена 


учебным материалом. Во-пер\-вых, в объеме заданий не предусматривается закрепление 


всех аспектов и нюансов знаний. Во-вто\-рых, количество заданий определяется исходя из 


целей контроля, а не из целей закрепления. Для закрепления нужно выполнение большого 


количества однотипных, но немного отличающихся задач. Для контроля достаточно 


одной задачи каждого типа.





Результат выполнения заданий обычно фиксируется, что может служить 


соответствующим отчетом, на основании которого учитель принимает решение о 


необходимости выполнения дополнительных заданий или о завершении темы. 


Правильной стратегией выполнения этого этапа информационной модели было бы 


проведение анализа неверных решений с целью определения тех аспектов нового знания, 


которые поняты учащимся неправильно. После чего следовало бы отправить учащегося на 


повторное изучение этих аспектов учебного материала, а затем снова предложить ему 


выполнение заданий.





Насколько сложившаяся практика соответствует детально изложенным элементам 


структуры учебного процесса?





Мотивационные материалы, если и существуют, то в самом зачаточном виде в 


предисловиях и введениях учебников. Разработчики считают, что внешняя 


привлекательность ЭОР и насыщенность его мультимедийным контентом повышают 


мотивацию. Следует, однако, отметить, что это второстепенный фактор. Учитель, когда 


понимает необходимость мотивационных действий, вынужден самостоятельно 


придумывать соответствующие материалы. Обратная связь о наличии мотивации 


отсутствует совершенно. Входная проверка наличия необходимых для работы с ЭОР 


знаний, как правило, отсутствует. Обратная связь о наличии восприятия осуществляется 


выборочно. Учитель не в состоянии проверить наличие восприятия всеми учениками 


класса точно так же, как не в состоянии убедиться в закреплении материала всеми. 


Учителю приходится самостоятельно собирать или изготавливать материалы для 


передачи учащимся на этапе восприятия. Учитель самостоятельно подбирает комплекты 


задач для своих учеников из различных источников. Ни один из источников не заявляет о 


полноте комплектов задач. Поскольку задача обеспечения полноты комплекта задач даже 


не ставится, то о полноте закрепления материала говорить не приходится.


{\looseness=1





}





Таким образом, очевидными проблемами процесса обучения являются:


\begin{itemize}


\item недостаточность или полное отсутствие мотивации;


\item отсутствие проверки наличия необходимых знаний перед передачей нового 


материала;


\item недостаточность или полное отсутствие обратной связи по восприятию;


\item недостаточность внимания к полноте закрепления.


\end{itemize}





Позволяют ли многочисленные ЭОР устранить или хотя 


бы уменьшить названные проблемы? Мотивационных материалов автору найти не 


удалось. Входной проверки знаний перед изложением материала также не обнаружилось. 


Соответственно, эти действия перекладываются на учителя. Точно так же, как и в 


отсутствие ЭОР, учителю приходится искать, подбирать или изготавливать необходимые 


ему материалы. О~полноте комплектов задач в ЭОР предлагается судить учителю. 


Подсказки, встроенные в задачи, объяснение сложностей, замена числовых данных при 


повторном предъявлении задач встречаются в отдельных ресурсах, но нигде не 


встречается все вместе, системно. Таким образом, те ЭОР, которые доступны в настоящее 


время, в по\-дав\-ля\-ющем большинстве не способны повысить эффективность учебного 


процесса, что и наблюдается в действительности.





\section{Какими должны быть электронные образовательные ресурсы}





Попробуем сформулировать необходимые требования к ЭОР, которые по\-мог\-ли бы решить 


указанные проблемы.





ЭОР, служащие для пробуждения мотивации у учащихся, должны содержать изложение 


места и значения предназначенного для восприятия материала в изучаемом предмете. По 


возможности он должен содержать общекультурное значение материала, а также значение 


этого материала для учащихся соответственно их возрасту и типовым увлечениям. ЭОР 


должен содержать демонстрацию преимуществ человека, владеющего изучаемым 


материалом.





Подача информации должна выстраиваться с использованием рекламных технологий. Для 


автоматизированной оценки возникшей мотивации должны быть разбросаны гиперссылки 


на более глубокое освещение материала, а также должны быть предложены пробные 


интересные и завлекательные задания, выполнить которые с легкостью может человек, 


освоивший предмет (тему), но представляющие известную трудность для 


неподготовленных учащихся. По наличию переходов по гиперссылкам и по попыткам 


выполнить задания можно судить о наличии мотивации. Соответствующий отчет с 


отображением уровня мотивации должен быть представлен учителю по завершении 


работы учащегося с ЭОР. На основании этого отчета учитель будет судить о том, 


достаточна ли возникшая мотивация учащегося и будет ли восприятие им материала 


успешным.





Разумным требованием было бы предварение мотивационным материалом каждого ЭОР, 


предназначенного для усвоения ка\-ко\-го-ли\-бо фрагмента учебной программы. Если 


мотивационный материал включен в состав ЭОР, то попытки учащихся перейти от 


демонстрации места и важности материала к непосредственному его изучению также 


могут послужить сигналом о достижении достаточного уровня мотивации.





В ЭОР, предназначенном для передачи информации по какой-либо теме учащимся, 


должны присутствовать следующие компоненты:


\begin{itemize}


\item проверка наличия необходимых для освоения темы знаний. В ЭОР должны быть 


предусмотрены краткие справ\-ки-на\-по\-ми\-на\-ния по забытому материалу для предъявления 


учащимся при непрохождении ими входной проверки по ка\-ким-ли\-бо заданиям. Должны 


быть также предусмотрены ссылки на неусвоенный материал в более тяжелых случаях;


\item сам материал по теме в объеме, необходимом и достаточном в соответствии с 


действующими нормативами. При этом материал должен, по возможности, подаваться в 


различном виде: в текстовом, голосовом, образном, схематическом и~т.\,п., чтобы 


учащиеся могли выбрать предпочтительный для их индивидуальных особенностей 


восприятия вид подачи материала;


\item расширенный и углубленный материал двух-трех уровней со ссылками, 


отправляющими к еще более углубленному изложению материала. Эти дополнительные 


уровни подачи материала, во-пер\-вых, способствуют более глубокому, а следовательно, 


более качественному изучению материала, а во-вто\-рых, служат маркерами фиксации 


восприятия материала;


\item механизм фиксации наличия восприятия материала. Механизм может опираться на 


контрольные вопросы, использование гиперссылок или других интерактивных действий, 


анализ времени пребывания на отдельных страницах и любые иные методы измерений.


\end{itemize}





Отчет по уровню восприятия должен предоставляться учителю по завершении работы 


учащегося с ЭОР. По этому отчету учитель делает заключение о достаточности уровня 


восприятия материала и возможности перехода к этапу закрепления.





Ресурс, несущий в себе набор задач на закрепление знания, а также формирование умений 


и навыков по теме, должен содержать:


\begin{itemize}


\item генератор упражнений и задач. Назначение генератора~--- создавать разные задачи 


для различных пользователей и для последовательного предъявления однотипных (но 


разных) задач одному пользователю;


\item статистический модуль сопоставления идентификатора учащегося и выданных и 


решенных задач, включая статистику попыток решения. ЭОР должен узнавать учащегося 


и вести список предъявленных ему задач, типов неудач с решениями, запрошенной 


помощи, использованных подсказок и~т.\,п.;


\item задачи всех типов, приобретение навыка решения которых необходимо для освоения 


последующих тем программы, а также для прохождения тестов и экзаменов;


\item модуль вариативности задач помимо задач нормальной сложности должен 


содержать как задачи более сложные, чем это необходимо по программе, так и задачи 


упрощенные для тех учащихся, которые не в состоянии сразу воспринять задачи 


требуемого уровня сложности;


\item модуль анализа встречающихся сложностей и выдачи соответствующих подсказок. 


В~случаях, когда подсказок оказывается недостаточно, должны быть предусмотрены 


ссылки на теоретический материал по теме для повторного его усвоения;


\item модуль выдачи результата закрепления материала по теме. Этот отчет 


свидетельствует о закреплении темы учащимся, и, вообще говоря, более не требуется 


никаких контрольных и тестовых заданий по этой теме.


\end{itemize}





Устроенный подобным образом ЭОР закрепления одновременно решает задачу контроля 


закрепления материала учащимся, поскольку завершающий отчет свидетельствует о том, 


что учащийся приобрел навык решения всех преду\-смот\-рен\-ных программой типов задач 


по теме.





\section{Заключение}





Как было показано выше, использование подавляющего большинства существующих в 


настоящее время ЭОР не только не приводит к повышению эффективности образования 


хоть по ка\-кой-то шкале, но и принципиально не может этого сделать. Дело в том, что 


существующие ЭОР за редким исключением практически не изменяют работы учителя. 


Что же касается работы (учебной деятельности) учащихся, то она также остается без 


изменений. Выявленные в статье компоненты информационного воздействия на 


учащегося, способные привести к повышению эффективности обучения, а именно: 


мотивационные материалы, обратная связь по мотивации, входная проверка знаний, 


обратная связь по уровню восприятия, полнота комплекта упражнений и задач и обратная 


связь по закреплению материала,~--- пока остаются вне пределов внимания авторов и 


создателей ЭОР. Тем не менее необходимость использования этих информационных 


компонентов непосредственно вытекает из теоретических положений педагогической 


психологии и педагогики.





\section{Тезаурус}





{\bfseries\textit{Восприятие}}~--- психологический процесс приема информации, 


получаемой при помощи органов чувств и отражения ее в сознании в виде образов 


субъективной картины мира. Восприятие является начальным этапом процесса познания~---


приобретения индивидуального опыта.





{\bfseries\textit{Мотивация}}~--- внутреннее побуждение к действию.





{\bfseries\textit{Научение}}~--- процесс и результат приобретения индивидуального опыта 


биологической системой (от простейших до человека как высшей формы ее организации в 


условиях Земли).





{\bfseries\textit{Обучение}}~--- целенаправленная, последовательная передача (трансляция) 


социокультурного (общественно-исторического) опыта другому человеку в специально 


созданных условиях.





{\bfseries\textit{Учение}}~--- научение человека в результате целенаправленного, 


сознательного присвоения им передаваемого (транслируемого) ему социокультурного 


(обще\-ст\-вен\-но-ис\-то\-ри\-че\-ско\-го) опыта и формируемого на этой основе 


индивидуального опыта. Следовательно, учение рассматривается как разновидность 


научения.





{\small\frenchspacing


{%\baselineskip=10.8pt


\addcontentsline{toc}{section}{Литература}


\begin{thebibliography}{9}





\bibitem{1-fed}


\Au{Христочевский С.\,А.} Почему использование ИКТ в образовании до сих пор не 


приводит к повышению качества образования~/\!\!/ Применение новых технологий в 


образовании: Мат-лы XXII Междунар. конф.~--- Троицк: Байтик, 2011. С.~208--209.


\bibitem{2-fed}


\Au{Федосеев А.\,А.} К~вопросу о недостаточном использовании электронных 


образовательных ресурсов~/\!\!/ Электронная Казань-2011: Мат-лы III Междунар. 


научно-прак\-тич. конф.~--- Казань: Юниверсум, 2011. С.~125--128.


\bibitem{3-fed}


\Au{Айсмонтас Б.\,Б.} Педагогическая психология: Электронный учебник. {\sf 


http://\linebreak www.ido.edu.ru/psychology/pedagogical\_psychology/index.html}.


\bibitem{4-fed}


\Au{Габай Т.\,В.} Педагогическая психология: Уч.\ пособие для студентов 


вузов.~--- 2-е изд., испр.~--- М.: Академия, 2005. 240~c.


\bibitem{5-fed}


\Au{Лернер И.\,Я.} Теория современного процесса обучения, ее значение для 


практики~/\!\!/ Советская педагогика, 1989. №\,11. С.~10--17.


\bibitem{6-fed}


Дидактика средней школы: Некоторые проблемы современной дидактики. {\sf 


http://\linebreak didaktica.ru}.


\bibitem{7-fed}


\Au{Бордовская Н.\,В., Реан А.\,А.} Педагогика: Уч-к для вузов.~--- СПб.: Питер, 


2000. 304~с.





    \label{end\stat}


\bibitem{8-fed}


\Au{Коменский Я.\,А., Локк Д., Руссо Ж.-Ж., Песталоцци~И.\,Г.} Педагогическое 


наследие.~--- М.: Педагогика, 1989. 416~с.


 \end{thebibliography}


}


}


